\section{Η Προ-\ENG{Standardization} Εποχή (1980-1993)}
\label{sec:pre_standardization}

Πριν από την εμφάνιση των πρώτων προτύπων \ENG{E-Learning} στα μέσα της δεκαετίας του 1990, η ανάπτυξη και διανομή εκπαιδευτικού περιεχομένου βασιζόταν σε αποσπασματικές, ιδιόκτητες λύσεις. Η περίοδος 1980-1993 χαρακτηρίζεται ως η «προ-\ENG{standardization} εποχή» και αποτελεί το πλαίσιο μέσα στο οποίο αναδύθηκε η ανάγκη για κοινά πρότυπα.

Κατά τη δεκαετία του 1980, τα \textit{\ENG{Computer-Based} \ENG{Training}} (\ENG{CBT}) συστήματα αποτελούσαν την κυρίαρχη μορφή ηλεκτρονικής εκπαίδευσης. Αυτά τα συστήματα λειτουργούσαν κυρίως σε αυτόνομους υπολογιστές (\ENG{standalone} \ENG{PCs}) και το περιεχόμενο διανεμόταν μέσω φυσικών μέσων αποθήκευσης, όπως δισκέτες 5.25" και 3.5", και αργότερα μέσω \ENG{CD-ROMs.} Κάθε \ENG{CBT} πρόγραμμα σχεδιαζόταν για συγκεκριμένο λειτουργικό σύστημα (π.χ. \ENG{MS-DOS}, \ENG{Apple} \ENG{II}, \ENG{Commodore}) και συχνά περιλάμβανε ενσωματωμένο κώδικα παρακολούθησης (\textit{\ENG{tracking}}) που αποθήκευε τα δεδομένα προόδου του χρήστη σε τοπικά αρχεία.

Η έλλειψη δικτυακής συνδεσιμότητας περιόριζε σημαντικά τη δυνατότητα κεντρικής διαχείρισης των εκπαιδευτικών προγραμμάτων. Οι διαχειριστές εκπαίδευσης δεν είχαν πρόσβαση σε πραγματικού χρόνου δεδομένα προόδου, και η συλλογή αποτελεσμάτων απαιτούσε χειροκίνητη εξαγωγή και συγκέντρωση αρχείων από κάθε σταθμό εργασίας. Επιπλέον, οι μορφές δεδομένων (\textit{\ENG{data} \ENG{formats}}) που χρησιμοποιούνταν ήταν εντελώς ιδιόκτητες και ασυμβίβαστες μεταξύ προμηθευτών, καθιστώντας αδύνατη την ενοποίηση στατιστικών στοιχείων από διαφορετικά \ENG{CBT} προγράμματα.

Στα τέλη της δεκαετίας του 1980 και στις αρχές της δεκαετίας του 1990, με την εξάπλωση των τοπικών δικτύων (\textit{\ENG{Local} \ENG{Area} \ENG{Networks}, \ENG{LANs}}), εμφανίστηκαν τα πρώτα πρωτόγονα συστήματα διαχείρισης μάθησης. Τα συστήματα αυτά επέτρεπαν την κεντρική αποθήκευση εκπαιδευτικού περιεχομένου και την παρακολούθηση της προόδου των εκπαιδευομένων μέσω δικτύου. Ωστόσο, η έλλειψη κοινών προδιαγραφών σήμαινε ότι κάθε \ENG{LMS} λειτουργούσε ως ένα «κλειστό οικοσύστημα», στο οποίο το εκπαιδευτικό υλικό ήταν στενά συνδεδεμένο με τη συγκεκριμένη πλατφόρμα. Ως αποτέλεσμα, ένα μάθημα που είχε αναπτυχθεί για ένα συγκεκριμένο \ENG{LMS} δεν μπορούσε εύκολα να μεταφερθεί σε άλλο χωρίς σημαντική επανεργασία.

Στην προ-\ENG{standardization} εποχή, οι οργανισμοί αντιμετώπιζαν μια σειρά από πρακτικά προβλήματα. Το κόστος ανάπτυξης εκπαιδευτικού περιεχομένου ήταν ιδιαίτερα υψηλό, καθώς κάθε νέο \ENG{LMS} απαιτούσε τη δημιουργία νέου υλικού ή την εκτεταμένη προσαρμογή του υπάρχοντος. Παράλληλα, οι οργανισμοί εξαρτώνταν έντονα από συγκεκριμένους προμηθευτές λογισμικού, γεγονός που περιόριζε τη δυνατότητα μετάβασης σε άλλες λύσεις. Επιπλέον, η απουσία κοινών προτύπων μέτρησης και αναφορών καθιστούσε δύσκολη τη σύγκριση και αξιολόγηση της αποτελεσματικότητας διαφορετικών εκπαιδευτικών προγραμμάτων.

Τα προβλήματα αυτά επηρέαζαν επίσης την κλιμάκωση των εκπαιδευτικών πρωτοβουλιών. Η επέκταση ενός προγράμματος κατάρτισης σε νέα τμήματα ή θυγατρικές εταιρείες συχνά απαιτούσε την επανάληψη μεγάλου μέρους της διαδικασίας ανάπτυξης. Τέλος, η ανταλλαγή εκπαιδευτικού υλικού μεταξύ διαφορετικών οργανισμών ή συστημάτων ήταν πρακτικά αδύνατη, γεγονός που περιόριζε σημαντικά τη συνεργασία και την επαναχρησιμοποίηση πόρων.

Η αναγνώριση αυτών των προβλημάτων, ιδίως από βιομηχανίες με υψηλές απαιτήσεις εκπαίδευσης (όπως η αεροπορία και η άμυνα), οδήγησε στην έναρξη συλλογικών προσπαθειών για τη δημιουργία κοινών προδιαγραφών. Η πρώτη τέτοια πρωτοβουλία ήταν η ίδρυση της \textit{\ENG{Aviation} \ENG{Industry} \ENG{Computer-Based} \ENG{Training} \ENG{Committee}} (\ENG{AICC}) το 1988.

Η προ-\ENG{standardization} εποχή αποτελεί σημαντικό ιστορικό στάδιο, καθώς καθιέρωσε τις βασικές έννοιες του \ENG{E-Learning} (παρακολούθηση προόδου, διαδραστικότητα, αξιολόγηση γνώσεων) και ανέδειξε τις τεχνικές και οικονομικές προκλήσεις που θα αντιμετώπιζαν τα μελλοντικά πρότυπα.

% \ENG{Figure} \ENG{showing} \ENG{pre-standardization} \ENG{problems}
\begin{figure}[H]
  \centering
  % FIGURE 2.2: Pre-Standardization Era Problems (Alternative Visualization)
\begin{chapterfigurebox}
\begin{tikzpicture}[
  node distance=2.5cm,
  every node/.style={font=\small},
  system/.style={draw, rounded corners, minimum width=3.2cm, minimum height=1.3cm, align=center, fill=problemred!15},
  content/.style={draw, rounded corners, minimum width=3cm, minimum height=1.2cm, align=center, fill=blue!15},
  problem/.style={draw, dashed, rounded corners, minimum width=3cm, minimum height=1cm, align=center, fill=orange!20},
  arrow/.style={->, thick}
]

% Content nodes
\node[content] (course1) {Μάθημα A\\\ENG{CBT Content}};
\node[content, below of=course1] (course2) {Μάθημα B\\\ENG{CBT Content}};
\node[content, below of=course2] (course3) {Μάθημα C\\\ENG{CBT Content}};

% LMS nodes
\node[system, right=4cm of course1] (lms1) {\ENG{LMS} 1\\Ιδιόκτητο \ENG{API}};
\node[system, right=4cm of course2] (lms2) {\ENG{LMS} 2\\Ιδιόκτητο \ENG{Format}};
\node[system, right=4cm of course3] (lms3) {\ENG{LMS} 3\\Διαφορετικό \ENG{Tracking}};

% Arrows
\draw[arrow] (course1) -- (lms1);
\draw[arrow] (course2) -- (lms2);
\draw[arrow] (course3) -- (lms3);

% Cross incompatibility
\draw[arrow, dashed] (course1.east) to[bend left=20] (lms2.west);
\draw[arrow, dashed] (course2.east) to[bend left=20] (lms3.west);
\draw[arrow, dashed] (course3.east) to[bend left=20] (lms1.west);

% Problems at bottom
\node[problem, below=2.2cm of course3, xshift=2cm] (p1) {Υψηλό κόστος\\επαναανάπτυξης};
\node[problem, right=2.5cm of p1] (p2) {Έλλειψη\\διαλειτουργικότητας};
\node[problem, right=2.5cm of p2] (p3) {Εξάρτηση\\από προμηθευτή};

% Title node
\node[font=\footnotesize, align=center, above=1cm of lms1] {
\textbf{Αποσπασματικό οικοσύστημα \ENG{CBT}/\ENG{LMS} πριν την τυποποίηση}
};

\end{tikzpicture}
\end{chapterfigurebox}
  \caption{Προβλήματα της Προ-\ENG{Standardization} Εποχής (1980-1993)}
  \label{fig:pre_standard_problems}
\end{figure}
